\section{Introduction}

A central problem in four-dimensional quantum topology is to construct a
tractable field-theoretic invariant that detects exotic smooth structure.
For premodular fusion categories, the Crane--Yetter theory
reduces to classical invariants and does not detect exotic smooth structure
\cite{Reutter2023}.  On the other hand, gauge-theoretic invariants have been
successful in detecting smooth structure, and second homology is part of
their input: Donaldson invariants admit insertions from $H_2$
\cite{Donaldson1990}, while Seiberg--Witten invariants are organized by
$\operatorname{Spin}^c$ structures and their classes in second cohomology
\cite{Witten1994}.  A natural idea is therefore to enrich the Crane--Yetter
theory by second-homology data and ask whether the resulting invariant can
detect exotic smoothness.

Yetter proposed precisely such a refinement in a sketch note in 1993.  For a
premodular fusion category $\cC$,\footnote{Yetter's terminology was an
``artinian semisimple tortile category.''  A tortile category is what is now
called a ribbon category.  With the finite set of simple isomorphism classes,
finite semisimplicity, rigidity, and simple tensor unit assumed here, this is
precisely a ribbon fusion category, equivalently a premodular fusion category.}
his sketch used the abelian coefficient group
\[
  A_{\cC}:=\operatorname{Aut}_{\otimes}(\operatorname{Id}_{\cC})
  \cong\operatorname{Hom}\bigl(U(\cC),\mathbb C^\times\bigr),
\]
    where $U(\cC)$ is the universal grading group
    \cite[Definitions~1 and~3 and Proposition~4]{Yetter1993}.  To avoid the
notation clash in that note, we reserve $U(\cC)$ for the universal grading
    group and $A_{\cC}$ for the coefficient group.  Yetter proposes homological
    twisting with coefficients in $A_{\cC}$, but the four-dimensional state sum
    actually displayed in Definition~6 and proved invariant in Theorem~8 is
    written for $H_2(M;\mathbb Z_2)$.  We instead fix
    $\eta\in A_{\cC}$ and use the coefficient homomorphism
    $\mathbb Z\to A_{\cC}$, $n\mapsto\eta^n$.  This is a cyclic integral
specialization of the general proposal, not a direct restatement of Yetter's
Theorem~8.  Yetter's note does not provide a Kirby-diagram handle formula, nor
does it determine whether the refinement detects exotic smooth structure.  We
address these two questions for the four-dimensional Crane--Yetter theory; we
do not consider the twisted Turaev--Viro theory.

The construction has three steps.  Fix a premodular fusion category $\cC$ over
$\mathbb C$ and a monoidal natural automorphism
\[
  \eta\colon\operatorname{Id}_{\cC}\Longrightarrow
  \operatorname{Id}_{\cC}.
\]
First, choose a marked Kirby presentation of $(X,\alpha)$.  Standard
\(0\)--\(1\) and \(3\)--\(4\) cancellations reduce it to a marked one-ended
presentation $(\mathcal H,\beta)$, where
$\beta\in H_2(|\mathcal H|;\mathbb Z)$ maps to $\alpha$.  At the evaluator
stage, Poincar\'e duality and the cellular comparison isomorphism turn
$\beta$ into a handle-cohomology class; choose a cellular $2$-cocycle
$\omega=(\omega_k)$ representing that class.  Second, decorate each framed $2$-handle by the
Kirby colour and the character coupon $(\eta^{\omega_k})$, and let every dotted
$1$-handle project the ribbons passing through it onto their tensor-unit
channel.  This gives the finite handle formula
\[
  \CY_{\cC,\eta}(\mathcal H,\omega)
  =\operatorname{Dim}(\cC)^{1-h_3(\mathcal H)}
    \sum_{\lambda}
      \left(\prod_kd_{\lambda_k}\right)
      \bigl\langle L^\eta_{\lambda,\omega}(\mathcal H)\bigr\rangle.
\]
Third, check the Kirby--Cerf moves together with their induced transport of
the handle cocycle.  The projector absorbs coboundary character coupons and
dotted-handle slides; monoidality of $\eta$ controls $2/2$ slides; and the
local cancellation factors match the global-dimension normalization.
Marked Kirby--Cerf completeness then produces an invariant
$\CY_{\cC,\eta}\bigl([(X,\alpha)]\bigr)$ of oriented diffeomorphism classes of marked
4-manifolds.  When $\alpha=0$, the character coupons are trivial and the
formula recovers the untwisted Crane--Yetter handle evaluation of
\cite[Definition~5.11]{Tham2021}.  We expect the construction to agree with
Yetter's four-dimensional state sum on their common domain, but that
identification is not needed for, and is not claimed as part of, our results.

The motivating question has a negative answer in the simply connected case:
$\CY_{\cC,\eta}\bigl([(X,\alpha)]\bigr)$ depends only on the
marked integral intersection lattice
\[
  \bigl(H_2(X;\mathbb Z),Q_X,\alpha\bigr),
\]
namely the integral second-homology group, its intersection pairing, and the
distinguished class $\alpha$, together with categorical constants.  For an odd
intersection form, the normalized value is a product of shifted positive and
negative Gauss sums evaluated on the diagonal coordinates of $\alpha$
(Proposition~\ref{prop:odd-formula}).  For an even intersection form, it is an
explicit stabilized quotient of the values on $K3$, reversed $K3$, and
hyperbolic blocks (Proposition~\ref{prop:even-formula}).  If $\cC$ is modular,
these two descriptions reduce to a single closed formula
(Proposition~\ref{prop:modular-formula}).  On setting the character twist to
be trivial, the odd and even formulas specialize respectively to the non-spin
and spin formulas of Reutter's Corollary~3.9, equations (5) and (6), for the
Crane--Yetter theory \cite{Reutter2023}.
Consequently this most direct semisimple homological refinement of the
Yetter--Crane--Yetter construction does not furnish an exotic-smoothness
witness for simply connected closed 4-manifolds
(Corollary~\ref{cor:no-simply-connected-exotic-detection}).

Reutter's untwisted obstruction does not require simple connectivity.  If two
closed oriented smooth 4-manifolds are orientation-preservingly
homeomorphic, an unmarked stable diffeomorphism is enough: the stabilization
factor is nonzero and can therefore be cancelled \cite{Reutter2023}.  The
marked theory requires strictly more.  Given a homeomorphism
$g\colon X\to Y$ and a class $\alpha\in H_2(X;\mathbb Z)$, the cancellation
argument requires a stable diffeomorphism $F$ satisfying
\[
  F_*(\alpha\oplus0)=(g_*\alpha)\oplus0.
\]
An arbitrary stable diffeomorphism need not satisfy this condition and may mix
$\alpha$ with the new hyperbolic summands.

In the simply connected setting, Wall's stable classification and the
one-stable realization theorem of Ruberman--Strle provide precisely the
required control \cite{Wall1964,RubermanStrle2022}.  This is where simple
connectivity enters our argument.  No corresponding marked realization
theorem is assumed here for general fundamental groups.  Consequently our
result rules out exotic-smoothness detection in the simply connected case,
while leaving the non-simply-connected marked case unresolved.

As in Reutter's calculation of semisimple four-dimensional field theories,
stabilization turns this invariant into an effective computational tool
\cite[Theorem~A and Corollary~3.9]{Reutter2023}.  Put
$H=S^2\times S^2$, write
$\widehat{\CY}=\operatorname{Dim}(\cC)^{-1}\CY$, and set
$h=\widehat{\CY}_{\cC,\eta}\bigl([(H,0)]\bigr)$.  Wall's theorem says that simply connected
4-manifolds with isometric intersection forms become diffeomorphic after
taking connected sums with copies of $H$ \cite[Theorems~2--3]{Wall1964}.
Because our invariant is marked by a class in $H_2$, we also use
Ruberman--Strle's one-stable realization theorem to arrange that the stable
diffeomorphism carries the chosen class to the chosen class
\cite[Theorem~2.1]{RubermanStrle2022}.  The connected-sum formula and the
direct calculation $h\neq0$ then imply that a marked stable diffeomorphism
\[
  (X\mathbin{\#}rH,\alpha\oplus0)
  \cong (M\mathbin{\#}sH,\beta\oplus0)
\]
determines the unstabilized value by
\[
  \widehat{\CY}_{\cC,\eta}\bigl([(X,\alpha)]\bigr)
  =\widehat{\CY}_{\cC,\eta}\bigl([(M,\beta)]\bigr)h^{s-r}.
\]
Thus stabilization is useful here not merely because it identifies more
manifolds, but because its nonzero contribution can be divided out.

This is especially effective for simply connected 4-manifolds.  Their
oriented homeomorphism types are controlled by their unimodular intersection
forms, together with the Kirby--Siebenmann invariant, which vanishes in the
smooth case \cite{Freedman1982}.  The odd indefinite forms are diagonal
\cite[Theorem~2.3.1]{PetersSterk2024}, while Donaldson's theorem gives the same
conclusion for smooth definite forms \cite[Theorem~A]{Donaldson1983}; their
standard models are connected sums of $\mathbb{CP}^2$ and
$\overline{\mathbb{CP}}^{\,2}$.  Even forms become, after hyperbolic
stabilization, connected sums of $K3$, $\overline{K3}$, and $H$ blocks, as in
the proof of \cite[Corollary~3.9]{Reutter2023}.  Consequently the preceding
stable calculation reduces the value on an arbitrary simply connected marked
manifold to values on these few standard pieces.  Connected-sum
multiplicativity and stabilization then give the odd and even closed formulas
in Propositions~\ref{prop:odd-formula} and
\ref{prop:even-formula}.  If $\cC$ is modular, the character--object
correspondence and the Gauss-sum identities of
\cite[Lemmas~8.15.3 and~8.22.9, Corollary~8.15.5, and
Proposition~8.15.4]{EGNO2015} reduce both formulas further to
Proposition~\ref{prop:modular-formula}.
